\scanpage{30}

\uline{\uline{III} \deleted{\unsure{1}}0}

\uline{Introduction}

Dans le par. 4, nous avons vu \deleted{que} lorsque $X$ est propre sur $Y$ \deleted{\unsure{un anneau local complet}} \inserted{\deleted{$I$-adique $A$}} \inserted{loc. noeth,} \deleted{et si on \unsure{connaît|considère}} et si $Y'$ est une partie fermée de $Y$, $X'=f^{-1}(Y')$, \deleted{la} \inserted{une propriété de} commutation des foncteurs \deleted{$R^i f$ et $R^i\hat f$ avec} ``images directes supérieures'' aux foncteurs complétion \unsure{Au long|(le long|le long} de $Y'$, resp $X'$]. Lorsque $Y$ était le spectre d'un anneau \deleted{local} noeth $I$-adique complet, et $Y'=V(A/I)$, ces résultats, complétés par un th d'existence (par 5), s'expriment \deleted{en disant} \inserted{en disant} \deleted{par} \inserted{que $F'\mathbin{-}F'_{/Y'}$ est une} équivalence de la catégorie des \deleted{\unsure{faisceaux}} \inserted{Modules} coh. sur $X'$ et de la catégorie des pro-Modules coh. sur $X'_{/Y'}$, \uline{compatible} avec la formation des modules de cohomologie.

Dans ce par, nous \deleted{généralisons} \inserted{montrons} que ces résultats peuvent se généraliser à certaines situations où $f:X'\to X$ n'est plus supposé propre [p. ex. $f$ une imm. ouverte \unsure{---|…}] \deleted{\unsure{ici}} la seule différence essentielle est que \deleted{\unsure{pr}} ici on \unsure{va trouver}
% slanted
\inserted{\unsure{dans l'énoncé des|des raisons ou} conditions \unsure{qui|que} pour des raisons \unsure{à voir|sont}, \unsure{par|pour} \unsure{nécessité|v. commentaire}.}
\deleted{qui s'expriment \unsure{d'une manière à préciser. Ceci|De manière à préciser ceci} \unsure{la vérité|le résultat} \unsure{essentiel est encore fourni} \unsure{un Th.}}
La méthode sera \inserted{\unsure{d'ailleurs toute parallèle à}} celle des par \deleted{\unsure{5|4}}4 et 5, \inserted{(cf N\textsuperscript{o}2)} en utilisant \deleted{des} \inserted{un} \uline{th. de finitude} convenable \inserted{(N\textsuperscript{o}1)} \deleted{\unsure{Ces de Ces derniers font}} \deleted{\unsure{faire intervenir}} faisant intervenir des propriétés \deleted{s'exprime :} \deleted{\unsure{intervenir reposent sur|intervenir reposant sur}} locales des faisceaux envisagés, qui s'expriment :

% archive sheet number: 29
